\section{Results}
\subsection{Status-linked differences in the frequency of culture-talk}
Table~\ref{tab:talk-tab1} shows the point estimates and associated $t$-statistics corresponding to the predictors of expected self-reported frequency of culture talk across all five domains. Hypothesis 1 predicts that a person's socio-economic status should predict culture talk even after adjusting for participation frequency. The results are consistent with this hypothesis. Looking at the effect of education, I find that there exist strong levels of status-linked heterogeneity in culture-talk even after holding constant self-reported consumption behavior. Persons who have attained advanced educational qualifications are more likely to engage in culture talk featuring popular culture and the arts as content than persons who have not, within levels of self-reported participation frequency. The results are also consistent with the claim that when it comes to culture talk, education is the more relevant status characteristic (in comparison to income). In contrast to education, income is only a statistically significant predictor of culture talk in three out of the five models, and the substantive strength of the effect is comparatively weaker.

\subsection{Heterogeneity in the Education Effect Levels of Participation}
How does the additive model compare to the multiplicative specification?  Figure~\ref{fig:educgraph} shows the effect of education on the marginal probability of engaging in culture ``very often'' for respondents partitioned according to the levels of self-reported consumption in that domain. The strength of the multiplicative interaction between education and consumption can be gauged by the extent to which we observe any trends in the marginal effect estimate. The estimates are obtained from ordered probit models specified according to equation~\ref{eq:eq2}. 

Hypothesis 2 predicts that the education effect should be more pronounced among culturally active persons.  The results provide fairly strong levels of support for this hypothesis. Upward trends in the effect of education on inducing frequent levels of culture talk are clear in film book-reading, and to a lesser extent music.  They are much weaker in the case of dramatic performances. In the case of the visual arts, I simply find a categorical distinction between the most active participants and everybody else, with education-based differentiation in culture talk being strong for the former group and non-existent for the latter.  In all, these results provide fairly broad support for the hypothesis that status linked differences in culture talk are much stronger for the most active culture consumers.

\subsection{Heterogeneity in the Socio-Economic Status Effect Across Cultural Domains}
Hypothesis 3 predicts that status-based differentiation should be stronger for the traditional high-status cultural domains, while hypothesis 4 predicts that we should find status-based differentiation in all domains.  A quick glance at the pattern of results shown in table~\ref{tab:talk-tab1} reveals that, when it comes to education, hypothesis 4 is on the right track. The education advantage in culture talk is comparably strong across all cultural domains. 

\subsection{Culture Talk Differences Across Commercial and Fine Arts Audiences}
Hypothesis 5 says that fine arts consumers should be as likely to engage in culture talk in commercial domains than commercial domain consumers. In order to test this hypothesis I proceed in two steps. First, I create two variables designed to tap the relative likelihood of each respondent to engage in the commercial and fine arts. These variables come from a principal factor analysis of the five consumption frequency items, in which I retain only the two main factors (those with an eigenvalue larger than 0.9). I use a promax rotation to allow the factors to be correlated with one another. The factor loadings for the first two principal factors are shown in~\ref{tab:fac-tab}. As we would expect reading, visiting museums and galleries and loading to the theater load on a common factor, while going to the movies and listening to the music load on a second factor. The regression-based scores from these two factors are taken as indicators of the relative likelihood to engage the fine and commercial arts respectively.

In a second step, I introduce these two variables as predictors of culture talk in an ordered probit model that adjusts for the same variables as those shown in table~\ref{tab:talk-tab1}. The coefficient estimates for these models are shown in table~\ref{tab:talk-tab2} (I omit the presentation of the other estimates for the sake of simplicity). To test the hypothesis that fine arts consumption is a better predictor of culture talk than commercial arts consumption across all domains, I estimate a constrained model that forces these two coefficients to be the same. I then use a likelihood ratio test to check whether there is a significant loss of model fit in comparison to the unconstrained model. The results provide strong support for hypothesis 5. Across all domains except for film, being a fine arts consumer is a better predictor of culture talk ($p<0.05$). 